% answers.tex —— 2014 年·上海卷·文科 · 参考答案与详细解析片段
% 由 tools/generate_projects.py 依来源总账自 2014/2014_数学真题_上海文_试题.md + 2014/2014_数学真题_上海文_解析正文.md + 2014/2014_数学真题_上海文_答案要点.md 生成，勿手改（改源 MD 后重跑）


\begin{answer}
\gksolhead{1}{}

\textbf{题目}：

函数 $y=1-2\cos^2(2x)$ 的最小正周期是\gkblank{}。

\Needspace{4\baselineskip}
\textbf{解析}

\Needspace{5\baselineskip}
由二倍角公式 $\cos^2(2x)=\dfrac{1+\cos4x}{2}$，得
\nopagebreak[4]
\[y=1-2\cos^2(2x)=1-(1+\cos4x)=-\cos4x,\]
\Needspace{5\baselineskip}
故最小正周期
\nopagebreak[4]
\[T=\frac{2\pi}{4}=\frac\pi2 .\]

\gkanswer{答案：$\dfrac\pi2$。由 $y=1-2\cos^2(2x)=-\cos4x$ 得 $T=\dfrac{2\pi}4=\dfrac\pi2$。}

\end{answer}
\begin{answer}
\gksolhead{2}{}

\textbf{题目}：

若复数 $z=1+2\mathrm i$，其中 $\mathrm i$ 是虚数单位，则 $\left(z+\dfrac1{\overline z}\right)\cdot\overline z=$\gkblank{}。

\Needspace{4\baselineskip}
\textbf{解析}

\Needspace{5\baselineskip}
由 $z=1+2\mathrm i$ 得 $\overline z=1-2\mathrm i$，$|z|^2=1^2+2^2=5$。利用 $\overline z\cdot\dfrac1{\overline z}=1$，
\nopagebreak[4]
\[\left(z+\frac1{\overline z}\right)\overline z=z\overline z+1=|z|^2+1=5+1=6 .\]

\gkanswer{答案：$6$。$\left(z+\dfrac1{\overline z}\right)\overline z=|z|^2+1=5+1=6$。}

\end{answer}
\begin{answer}
\gksolhead{3}{}

\textbf{题目}：

设常数 $a\in\mathbf R$，函数 $f(x)=|x-1|+|x^2-a|$，若 $f(2)=1$，则 $f(1)=$\gkblank{}。

\Needspace{4\baselineskip}
\textbf{解析}

\Needspace{5\baselineskip}
由 $f(2)=1$ 得
\nopagebreak[4]
\[f(2)=|2-1|+|2^2-a|=1+|4-a|=1,\]
\Needspace{5\baselineskip}
故 $|4-a|=0$，即 $a=4$。于是
\nopagebreak[4]
\[f(1)=|1-1|+|1^2-4|=0+3=3 .\]

\gkanswer{答案：$3$。$f(2)=1+|4-a|=1$ 得 $a=4$，故 $f(1)=|0|+|1-4|=3$。}

\end{answer}
\begin{answer}
\gksolhead{4}{}

\textbf{题目}：

若抛物线 $y^2=2px$ 的焦点与椭圆 $\dfrac{x^2}{9}+\dfrac{y^2}{5}=1$ 的右焦点重合，则该抛物线的准线方程为\gkblank{}。

\Needspace{4\baselineskip}
\textbf{解析}

\Needspace{5\baselineskip}
椭圆 $\dfrac{x^2}9+\dfrac{y^2}5=1$ 中 $a^2=9$，$b^2=5$，故
\nopagebreak[4]
\[c=\sqrt{a^2-b^2}=\sqrt{9-5}=2,\]
右焦点为 $(2,0)$。

\Needspace{5\baselineskip}
抛物线 $y^2=2px$ 的焦点为 $\left(\dfrac p2,0\right)$，由 $\dfrac p2=2$ 得 $p=4$，其准线为
\nopagebreak[4]
\[x=-\frac p2=-2 .\]

\gkanswer{答案：$x=-2$。椭圆右焦点为 $(2,0)$，故 $\dfrac p2=2$，$p=4$，准线 $x=-\dfrac p2=-2$。}

\end{answer}
\begin{answer}
\gksolhead{5}{}

\textbf{题目}：

某校高一、高二、高三分别有学生 $1600$ 名、$1200$ 名、$800$ 名。为了解该校高中学生的牙齿健康状况，按各年级的学生数进行分层抽样。若高三抽出 $20$ 名学生，则高一、高二共抽取的学生数为\gkblank{}。

\Needspace{4\baselineskip}
\textbf{解析}

\Needspace{5\baselineskip}
分层抽样按各年级人数比例抽取，抽样比为
\nopagebreak[4]
\[\frac{20}{800}=\frac1{40}.\]
\Needspace{5\baselineskip}
高一、高二共有 $1600+1200=2800$ 名学生，故共抽取
\nopagebreak[4]
\[\frac{2800}{40}=70\]
名。

\gkanswer{答案：$70$。抽样比为 $\dfrac{20}{800}=\dfrac1{40}$，故高一、高二共抽取 $\dfrac{1600+1200}{40}=70$ 名。}

\end{answer}
\begin{answer}
\gksolhead{6}{}

\textbf{题目}：

若实数 $x$，$y$ 满足 $xy=1$，则 $x^2+2y^2$ 的最小值为\gkblank{}。

\Needspace{4\baselineskip}
\textbf{解析}

\Needspace{5\baselineskip}
由 $xy=1$ 知 $x$，$y$ 同号且都不为零，于是
\nopagebreak[4]
\[x^2+2y^2\geqslant2\sqrt{x^2\cdot2y^2}=2\sqrt2\,|xy|=2\sqrt2 .\]
等号成立当且仅当 $x^2=2y^2$ 且 $xy=1$，例如 $x=2^{\frac14}$，$y=2^{-\frac14}$，此时 $xy=1$ 且 $x^2+2y^2=2\sqrt2$。故最小值为 $2\sqrt2$。

\gkanswer{答案：$2\sqrt2$。$x^2+2y^2\geqslant2\sqrt2|xy|=2\sqrt2$，$x^2=2y^2$ 且 $xy=1$ 时取等号。}

\end{answer}
\begin{answer}
\gksolhead{7}{}

\textbf{题目}：

在长方体中割去两个小长方体后的几何体的三视图如图，则切割掉的两个小长方体的体积之和等于\gkblank{}。
\nopagebreak[4]
\begin{center}% ============================================================
% f2014_shanghai_arts_q07_views.tex
% 2014 上海卷（文科）第 7 题：割去两个小长方体后的三视图
% 来源 PDF：普通高考/2014/2014上海文.pdf
% 原页：第 1 页（A3），印面 917，填空题第 7 题题干下方
% 原长方体 5（长）×3（深）×4（高）；中间高 2 的一层左右两侧
% 各割去一个 2×2×3 的小长方体，正中留下宽 1 的竖条。
% 正视图：宽 5、高 4 的“工”字形；侧视图：3×4 矩形，高 1、3 处
% 各有实线；俯视图：5×3 矩形，距左边 2、3 处各有一条贯通虚线。
% ============================================================
\begin{tikzpicture}[>=Stealth, line width=0.65pt, font=\small, scale=0.86,
  line cap=round, line join=round,
  dim/.style={<->, >=Stealth, thin, line width=0.45pt}]

  % ---------------- 正视图 ----------------
  \draw (0,0) -- (5,0) -- (5,1) -- (3,1) -- (3,3) -- (5,3) -- (5,4)
        -- (0,4) -- (0,3) -- (2,3) -- (2,1) -- (0,1) -- cycle;

  % 中间竖条宽 1 的尺寸线（延长线伸入上横带）
  \draw[thin] (2,3) -- (2,3.78);
  \draw[thin] (3,3) -- (3,3.78);
  \draw[->, thin, line width=0.45pt] (1.55,3.55) -- (1.96,3.55);
  \draw[->, thin, line width=0.45pt] (3.45,3.55) -- (3.04,3.55);
  \node at (2.5,3.55) {$1$};

  % 左侧三段高度 1、2、1
  \draw[dim] (-0.45,0) -- (-0.45,1);
  \draw[dim] (-0.45,1) -- (-0.45,3);
  \draw[dim] (-0.45,3) -- (-0.45,4);
  \foreach \y in {0,1,3,4} \draw[thin] (-0.62,\y) -- (-0.28,\y);
  \node[left] at (-0.62,0.5) {$1$};
  \node[left] at (-0.62,2) {$2$};
  \node[left] at (-0.62,3.5) {$1$};

  % 底边宽 5
  \draw[dim] (0,-0.5) -- (5,-0.5);
  \draw[thin] (0,-0.3) -- (0,-0.7);
  \draw[thin] (5,-0.3) -- (5,-0.7);
  \node[below] at (2.5,-0.5) {$5$};

  % ---------------- 侧视图 ----------------
  \draw (7,0) rectangle (10,4);
  \draw (7,1) -- (10,1);
  \draw (7,3) -- (10,3);
  \draw[dim] (7,-0.5) -- (10,-0.5);
  \draw[thin] (7,-0.3) -- (7,-0.7);
  \draw[thin] (10,-0.3) -- (10,-0.7);
  \node[below] at (8.5,-0.5) {$3$};

  % ---------------- 俯视图 ----------------
  \draw (0,-6) rectangle (5,-3);
  \draw[dashed] (2,-6) -- (2,-3);
  \draw[dashed] (3,-6) -- (3,-3);

  % 视图名称
  \node at (2.5,4.55) {正视图};
  \node at (8.5,4.55) {侧视图};
  \node at (2.5,-6.55) {俯视图};
\end{tikzpicture}
\end{center}
\nopagebreak[4]
\Needspace{4\baselineskip}
\textbf{解析}

由三视图读出原长方体的尺寸：正视图总宽 $5$、总高 $4=1+2+1$，上、下两条高为 $1$ 的横带贯通全宽，中间高为 $2$ 的一层只在正中留下宽 $1$ 的竖条；侧视图是宽 $3$、高 $4$ 的矩形，并在高 $1$、$3$ 处各有一条实线；俯视图是 $5\times3$ 的矩形，内部有距左边 $2$、$3$ 的两条贯通全深的虚线。

正视图缺口开在中间层的左右两侧（缺口处正视图既无轮廓也无横向实线），侧视图的两条实线正是缺口上、下边界，俯视图的两条虚线正是缺口的两个侧面，故两个缺口都贯通全深。于是中间层左右两侧各割去一个长 $2$、高 $2$、深 $3$ 的小长方体。

\Needspace{5\baselineskip}
每个小长方体体积为 $2\times2\times3=12$，故体积之和为
\nopagebreak[4]
\[2\times12=24 .\]
\Needspace{5\baselineskip}
检验：原长方体体积 $5\times3\times4=60$，余下部分体积
\nopagebreak[4]
\[2\times(5\times3\times1)+1\times3\times2=30+6=36,\]
而 $60-36=24$，与三视图一致。

\gkanswer{答案：$24$。三视图还原出中间层左右各割去一个 $2\times2\times3$ 的小长方体，体积之和为 $2\times12=24$。}

\end{answer}
\begin{answer}
\gksolhead{8}{}

\textbf{题目}：

若圆锥的侧面积是底面积的 $3$ 倍，则其母线与轴所成角的大小为\gkblank{}（结果用反三角函数值表示）。

\Needspace{4\baselineskip}
\textbf{解析}

\Needspace{5\baselineskip}
设底面半径为 $r$，母线长为 $l$，母线与轴所成角为 $\theta$。圆锥侧面积为 $\pi rl$，底面积为 $\pi r^2$，由题意
\nopagebreak[4]
\[\pi rl=3\pi r^2\iff l=3r .\]
\Needspace{5\baselineskip}
母线与轴、底面半径构成直角三角形，其中斜边为 $l$，角 $\theta$ 的对边为 $r$，故
\nopagebreak[4]
\[\sin\theta=\frac rl=\frac13,\]
从而 $\theta=\arcsin\dfrac13$。

\gkanswer{答案：$\arcsin\dfrac13$。由 $\pi rl=3\pi r^2$ 得 $l=3r$，母线与轴所成角 $\theta$ 满足 $\sin\theta=\dfrac rl=\dfrac13$。}

\end{answer}
\begin{answer}
\gksolhead{9}{}

\textbf{题目}：

\Needspace{5\baselineskip}
设
\nopagebreak[4]
\[f(x)=\begin{cases}
(x-a)^2, & x\leqslant0,\\
x+\dfrac1x+a, & x>0,
\end{cases}\]
若 $f(0)$ 是 $f(x)$ 的最小值，则 $a$ 的取值范围为\gkblank{}。

\Needspace{4\baselineskip}
\textbf{解析}

$f(0)=(0-a)^2=a^2$。

\Needspace{5\baselineskip}
当 $x\leqslant0$ 时，$f(x)=(x-a)^2$。要求 $f(x)\geqslant f(0)=a^2$ 对一切 $x<0$ 成立：
\nopagebreak[4]
\[(x-a)^2\geqslant a^2\iff x^2-2ax\geqslant0\iff x(x-2a)\geqslant0 .\]
\Needspace{5\baselineskip}
因 $x<0$，该式等价于 $x-2a\leqslant0$，即 $a\geqslant\dfrac x2$ 对一切 $x<0$ 成立，故
\nopagebreak[4]
\[a\geqslant0 .\]

\Needspace{5\baselineskip}
当 $x>0$ 时，由基本不等式
\nopagebreak[4]
\[f(x)=x+\frac1x+a\geqslant2+a,\]
\Needspace{5\baselineskip}
等号在 $x=1$ 时成立。要求 $a^2\leqslant2+a$，即
\nopagebreak[4]
\[a^2-a-2\leqslant0\iff-1\leqslant a\leqslant2 .\]

\Needspace{5\baselineskip}
两条件合并得
\nopagebreak[4]
\[0\leqslant a\leqslant2 .\]
当 $a\in[0,2]$ 时，$x\leqslant0$ 段最小值 $a^2$ 在 $x=0$ 取到，$x>0$ 段最小值 $2+a\geqslant a^2$，故 $f(0)$ 确为全局最小值。

\gkanswer{答案：$[0,2]$。$x\leqslant0$ 一侧要求 $a\geqslant0$；$x>0$ 一侧要求 $a^2\leqslant a+2$，即 $-1\leqslant a\leqslant2$。}

\end{answer}
\begin{answer}
\gksolhead{10}{}

\textbf{题目}：

设无穷等比数列 $\{a_n\}$ 的公比为 $q$，若 $a_1=\lim\limits_{n\to\infty}(a_3+a_4+\cdots+a_n)$，则 $q=$\gkblank{}。

\Needspace{4\baselineskip}
\textbf{解析}

\Needspace{5\baselineskip}
因 $\lim\limits_{n\to\infty}(a_3+a_4+\cdots+a_n)$ 存在，故 $|q|<1$ 且 $q\neq0$；又 $\{a_n\}$ 为等比数列，$a_1\neq0$。由 $a_k=a_1q^{k-1}$ 得
\nopagebreak[4]
\[\lim_{n\to\infty}(a_3+a_4+\cdots+a_n)=\frac{a_3}{1-q}=\frac{a_1q^2}{1-q}.\]
\Needspace{5\baselineskip}
由 $a_1=\dfrac{a_1q^2}{1-q}$ 约去 $a_1$ 得
\nopagebreak[4]
\[1=\frac{q^2}{1-q}\iff q^2+q-1=0,\]
\Needspace{5\baselineskip}
解得 $q=\dfrac{-1\pm\sqrt5}{2}$。结合 $|q|<1$ 得
\nopagebreak[4]
\[q=\frac{\sqrt5-1}{2}.\]

\gkanswer{答案：$\dfrac{\sqrt5-1}{2}$。由 $a_1=\dfrac{a_1q^2}{1-q}$ 得 $q^2+q-1=0$，结合 $|q|<1$ 取正根。}

\end{answer}
\begin{answer}
\gksolhead{11}{}

\textbf{题目}：

若 $f(x)=x^{\frac23}-x^{-\frac12}$，则满足 $f(x)<0$ 的 $x$ 的取值范围是\gkblank{}。

\Needspace{4\baselineskip}
\textbf{解析}

\Needspace{5\baselineskip}
由 $x^{-\frac12}$ 有意义得 $x>0$。不等式 $f(x)<0$ 即
\nopagebreak[4]
\[x^{\frac23}<x^{-\frac12},\]
\Needspace{5\baselineskip}
两边同乘 $x^{\frac12}>0$ 得
\nopagebreak[4]
\[x^{\frac23+\frac12}=x^{\frac76}<1,\]
故 $0<x<1$，取值范围为 $(0,1)$。

\gkanswer{答案：$(0,1)$。$x>0$，$x^{\frac23}<x^{-\frac12}\iff x^{\frac76}<1\iff0<x<1$。}

\end{answer}
\begin{answer}
\gksolhead{12}{}

\textbf{题目}：

方程 $\sin x+\sqrt3\cos x=1$ 在区间 $[0,2\pi]$ 上的所有解的和等于\gkblank{}。

\Needspace{4\baselineskip}
\textbf{解析}

\Needspace{5\baselineskip}
由辅助角公式
\nopagebreak[4]
\[\sin x+\sqrt3\cos x=2\sin\left(x+\frac\pi3\right).\]
\Needspace{5\baselineskip}
令 $t=x+\dfrac\pi3$，当 $x\in[0,2\pi]$ 时 $t\in\left[\dfrac\pi3,\dfrac{7\pi}3\right]$，原方程化为
\nopagebreak[4]
\[\sin t=\frac12 .\]
\Needspace{5\baselineskip}
在该区间上满足 $\sin t=\dfrac12$ 的角为
\nopagebreak[4]
\[t=\frac{5\pi}6,\qquad t=\frac{13\pi}6,\]
\Needspace{5\baselineskip}
对应
\nopagebreak[4]
\[x_1=\frac{5\pi}6-\frac\pi3=\frac\pi2,\qquad x_2=\frac{13\pi}6-\frac\pi3=\frac{11\pi}6 .\]
\Needspace{5\baselineskip}
（$t=\dfrac\pi6$ 不在区间内。）故所有解的和为
\nopagebreak[4]
\[x_1+x_2=\frac\pi2+\frac{11\pi}6=\frac{3\pi+11\pi}6=\frac{7\pi}3 .\]

\gkanswer{答案：$\dfrac{7\pi}3$。$2\sin\left(x+\dfrac\pi3\right)=1$ 在 $[0,2\pi]$ 上的解为 $\dfrac\pi2$，$\dfrac{11\pi}6$，其和为 $\dfrac{7\pi}3$。}

\end{answer}
\begin{answer}
\gksolhead{13}{}

\textbf{题目}：

为强化安全意识，某商场拟在未来的连续 $10$ 天中随机选择 $3$ 天进行紧急疏散演练，则选择的 $3$ 天恰好为连续 $3$ 天的概率是\gkblank{}（结果用最简分数表示）。

\Needspace{4\baselineskip}
\textbf{解析}

\Needspace{5\baselineskip}
从 $10$ 天中任选 $3$ 天，共有
\nopagebreak[4]
\[\mathrm C_{10}^3=\frac{10\times9\times8}{3\times2\times1}=120\]
种等可能的选法。

\Needspace{5\baselineskip}
$3$ 天恰好连续的情形为 $\{1,2,3\}$，$\{2,3,4\}$，$\cdots$，$\{8,9,10\}$，共 $8$ 种。故所求概率
\nopagebreak[4]
\[P=\frac8{120}=\frac1{15}.\]

\gkanswer{答案：$\dfrac1{15}$。连续三天共 $8$ 种，$\dfrac8{\mathrm C_{10}^3}=\dfrac8{120}=\dfrac1{15}$。}

\end{answer}
\begin{answer}
\gksolhead{14}{}

\textbf{题目}：

已知曲线 $C:x=-\sqrt{4-y^2}$，直线 $l:x=6$。若对于点 $A(m,0)$，存在 $C$ 上的点 $P$ 和 $l$ 上的点 $Q$ 使得 $\overrightarrow{AP}+\overrightarrow{AQ}=\vec{0}$，则 $m$ 的取值范围为\gkblank{}。

\Needspace{4\baselineskip}
\textbf{解析}

\Needspace{5\baselineskip}
$C:x=-\sqrt{4-y^2}$ 是圆 $x^2+y^2=4$ 在 $y$ 轴左侧的半圆（含端点），故 $P(x_P,y_P)$ 满足
\nopagebreak[4]
\[x_P\in[-2,0],\qquad x_P^2+y_P^2=4 .\]

\Needspace{5\baselineskip}
设 $Q(6,y_Q)$。由 $\overrightarrow{AP}+\overrightarrow{AQ}=\vec{0}$ 知 $A$ 是线段 $PQ$ 的中点，故
\nopagebreak[4]
\[m=\frac{x_P+6}{2},\qquad 0=\frac{y_P+y_Q}{2}.\]
\Needspace{5\baselineskip}
由 $x_P\in[-2,0]$ 得
\nopagebreak[4]
\[m\in\left[\frac{-2+6}2,\frac{0+6}2\right]=[2,3].\]
反之，任给 $m\in[2,3]$，取 $x_P=2m-6\in[-2,0]$，令 $y_P=\sqrt{4-x_P^2}$，$y_Q=-y_P$，则 $P\in C$，$Q(6,y_Q)\in l$ 且 $\overrightarrow{AP}+\overrightarrow{AQ}=\vec{0}$。故 $m$ 的取值范围为 $[2,3]$。

\gkanswer{答案：$[2,3]$。$A$ 为 $PQ$ 中点，$m=\dfrac{x_P+6}2$，而 $x_P\in[-2,0]$。}

\end{answer}
\begin{answer}
\gksolhead{15}{}

\textbf{题目}：

设 $a$，$b\in\mathbf R$，则“$a+b>4$”是“$a>2$ 且 $b>2$”的\mbox{（　　）}

（A）充分非必要条件；（B）必要非充分条件；（C）充要条件；（D）既不充分也不必要条件。

\Needspace{4\baselineskip}
\textbf{解析}

若 $a>2$ 且 $b>2$，则 $a+b>4$，故“$a+b>4$”是“$a>2$ 且 $b>2$”的必要条件。

反之，取 $a=1$，$b=4$，则 $a+b=5>4$，但 $a>2$ 不成立，故“$a+b>4$”不是充分条件。

所以“$a+b>4$”是“$a>2$ 且 $b>2$”的必要非充分条件。

选 B。

\gkanswer{答案：B。$a>2$ 且 $b>2$ 推出 $a+b>4$，反之 $a=1$，$b=4$ 为反例。}

\end{answer}
\begin{answer}
\gksolhead{16}{}

\textbf{题目}：

已知互异的复数 $a$，$b$ 满足 $ab\neq0$，集合 $\{a,b\}=\{a^2,b^2\}$，则 $a+b=$\mbox{（　　）}

（A）$2$；（B）$1$；（C）$0$；（D）$-1$。

\Needspace{4\baselineskip}
\textbf{解析}

由 $\{a,b\}=\{a^2,b^2\}$ 且 $a\neq b$，只有两种对应方式。

若 $a^2=a$ 且 $b^2=b$，则 $a$，$b\in\{0,1\}$，与 $ab\neq0$ 及 $a\neq b$ 矛盾。

\Needspace{5\baselineskip}
若 $a^2=b$ 且 $b^2=a$，则 $a^4=a$，即
\nopagebreak[4]
\[a(a-1)(a^2+a+1)=0 .\]
\Needspace{5\baselineskip}
由 $ab\neq0$ 得 $a\neq0$；由 $a\neq b$ 得 $a\neq1$（否则 $b=a^2=1=a$）。故 $a^2+a+1=0$，即 $a$ 是三次单位根 $\omega$（$\omega\neq1$，$\omega^3=1$），且 $b=a^2$。于是
\nopagebreak[4]
\[a+b=a+a^2=-1 .\]

选 D。

\gkanswer{答案：D。必有 $a^2=b$，$b^2=a$（$a^2=a$ 一支与 $ab\neq0$、$a\neq b$ 矛盾），得 $a^2+a+1=0$，故 $a+b=a+a^2=-1$。}

\end{answer}
\begin{answer}
\gksolhead{17}{}

\textbf{题目}：

如图，四个边长为 $1$ 的小正方形排成一个大正方形，$AB$ 是大正方形的一边，$P_i$（$i=1,2,\cdots,7$）是小正方形的其余顶点，则 $\overrightarrow{AB}\cdot\overrightarrow{AP_i}$（$i=1,2,\cdots,7$）的不同值的个数为\mbox{（　　）}
\nopagebreak[4]
\begin{center}% ============================================================
% f2014_shanghai_arts_q17_square.tex
% 2014 上海卷（文科）第 17 题：四个小正方形拼成的大正方形
% 来源 PDF：普通高考/2014/2014上海文.pdf
% 原页：第 1 页（A3），印面 917，选择题第 17 题题干下方
% 大正方形边长 2，内含两条中位线；九个格点中 A、B 为大正方形
% 左下、左上顶点（AB 为左边），其余七个格点标为 P1—P7：
%   P1 左边中点，P2 下边中点，P3 中心，P4 上边中点，
%   P5 右下顶点，P6 右边中点，P7 右上顶点。
% ============================================================
\begin{tikzpicture}[>=Stealth, line width=0.65pt, font=\small, scale=2.1,
  line cap=round, line join=round]
  % 大正方形及两条中位线
  \draw (0,0) rectangle (2,2);
  \draw (1,0) -- (1,2);
  \draw (0,1) -- (2,1);

  % 九个格点
  \fill (0,0) circle (0.75pt);
  \fill (0,1) circle (0.75pt);
  \fill (0,2) circle (0.75pt);
  \fill (1,0) circle (0.75pt);
  \fill (1,1) circle (0.75pt);
  \fill (1,2) circle (0.75pt);
  \fill (2,0) circle (0.75pt);
  \fill (2,1) circle (0.75pt);
  \fill (2,2) circle (0.75pt);

  % 顶点标注
  \node[below left=-1pt] at (0,0) {$A$};
  \node[above left=-1pt] at (0,2) {$B$};
  \node[left=-1pt] at (0,1) {$P_1$};
  \node[below=-1pt] at (1,0) {$P_2$};
  \node[above right=-2pt] at (1,1) {$P_3$};
  \node[above=-1pt] at (1,2) {$P_4$};
  \node[below right=-1pt] at (2,0) {$P_5$};
  \node[right=-1pt] at (2,1) {$P_6$};
  \node[above right=-1pt] at (2,2) {$P_7$};
\end{tikzpicture}
\end{center}
\nopagebreak[4]
（A）$7$；（B）$5$；（C）$3$；（D）$1$。

\Needspace{4\baselineskip}
\textbf{解析}

以大正方形的顶点为格点建立坐标系，取 $A(0,0)$，$B(0,2)$，则大正方形为 $[0,2]\times[0,2]$，九个格点为 $(i,j)$（$i,j\in\{0,1,2\}$）。$A$，$B$ 占去其中两个，其余七个点 $P_i$ 的坐标恰好遍历另外七组 $(i,j)$。

\Needspace{5\baselineskip}
于是
\nopagebreak[4]
\[\overrightarrow{AB}=(0,2),\qquad \overrightarrow{AP_i}=(i,j),\]
\nopagebreak[4]
\[\overrightarrow{AB}\cdot\overrightarrow{AP_i}=0\cdot i+2\cdot j=2j .\]
\Needspace{5\baselineskip}
七个点中 $j$ 取遍 $0$，$1$，$2$，故点积只取
\nopagebreak[4]
\[0,\quad2,\quad4\]
三个不同值。不同值的个数为 $3$。

选 C。

\gkanswer{答案：C。取 $A(0,0)$，$B(0,2)$，七个格点的纵坐标为 $0,1,2$，$\overrightarrow{AB}\cdot\overrightarrow{AP_i}=2y_i$ 只有 $0$，$2$，$4$ 三个值。}

\end{answer}
\begin{answer}
\gksolhead{18}{}

\textbf{题目}：

\Needspace{5\baselineskip}
已知 $P_1(a_1,b_1)$ 与 $P_2(a_2,b_2)$ 是直线 $y=kx+1$（$k$ 为常数）上两个不同的点，则关于 $x$ 和 $y$ 的方程组
\nopagebreak[4]
\[\begin{cases}
a_1x+b_1y=1,\\
a_2x+b_2y=1
\end{cases}\]
的解的情况是\mbox{（　　）}

（A）无论 $k$，$P_1$，$P_2$ 如何，总是无解；（B）无论 $k$，$P_1$，$P_2$ 如何，总有唯一解；（C）存在 $k$，$P_1$，$P_2$，使之恰有两解；（D）存在 $k$，$P_1$，$P_2$，使之有无穷多解。

\Needspace{4\baselineskip}
\textbf{解析}

\Needspace{5\baselineskip}
由 $P_i(a_i,b_i)$ 在直线 $y=kx+1$ 上得
\nopagebreak[4]
\[b_i=ka_i+1\qquad(i=1,2).\]
\Needspace{5\baselineskip}
方程组的系数行列式
\nopagebreak[4]
\[D=\begin{vmatrix}a_1&b_1\\a_2&b_2\end{vmatrix}=a_1b_2-a_2b_1=a_1(ka_2+1)-a_2(ka_1+1)=a_1-a_2 .\]
若 $a_1=a_2$，则 $b_1=ka_1+1=ka_2+1=b_2$，即 $P_1$，$P_2$ 重合，与题设矛盾。故 $D=a_1-a_2\neq0$，方程组恒有唯一解。

\Needspace{5\baselineskip}
事实上取 $x=-k$，$y=1$，则
\nopagebreak[4]
\[a_ix+b_iy=-ka_i+(ka_i+1)=1\qquad(i=1,2),\]
故解就是 $(-k,1)$，无论 $k$，$P_1$，$P_2$ 如何都是唯一解。

选 B。

\gkanswer{答案：B。由 $b_i=ka_i+1$ 得 $D=a_1-a_2\neq0$，且 $(-k,1)$ 恒为解。}

\end{answer}
\begin{answer}
\gksolhead{19}{}

\textbf{题目}：

底面边长为 $2$ 的正三棱锥 $P\text{-}ABC$，其表面展开图是三角形 $P_1P_2P_3$，如图，求 $\triangle P_1P_2P_3$ 的各边长及此三棱锥的体积 $V$。
\nopagebreak[4]
\begin{center}% ============================================================
% f2014_shanghai_arts_q19_net.tex
% 2014 上海卷（文科）第 19 题：正三棱锥 P-ABC 的表面展开图
% 来源 PDF：普通高考/2014/2014上海文.pdf
% 原页：第 1 页（A3），印面 917，解答题第 19 题题干下方（中栏下部）
% 展开图为三角形 P1P2P3：A、C 分别为 P1P3、P3P2 的中点，
% B 为 P1P2 的中点；△ABC 为底面（实线），大三角形三边为
% 展开图外边界（虚线）。
% ============================================================
\begin{tikzpicture}[>=Stealth, line width=0.65pt, font=\small, scale=0.8,
  line cap=round, line join=round]
  \coordinate (P1) at (0,0);
  \coordinate (P2) at (4,0);
  \coordinate (P3) at (2,3.4641);
  \coordinate (A)  at (1,1.73205);
  \coordinate (B)  at (2,0);
  \coordinate (C)  at (3,1.73205);

  % 展开图外边界（虚线）
  \draw[dashed] (P1) -- (P2) -- (P3) -- cycle;

  % 底面三角形（折痕，实线）
  \draw (A) -- (B) -- (C) -- cycle;

  % 顶点标注
  \fill (P1) circle (0.85pt) node[below left=-1pt] {$P_1$};
  \fill (P2) circle (0.85pt) node[below right=-1pt] {$P_2$};
  \fill (P3) circle (0.85pt) node[above=-1pt] {$P_3$};
  \fill (A) circle (0.85pt) node[above left=-1pt] {$A$};
  \fill (B) circle (0.85pt) node[below=-1pt] {$B$};
  \fill (C) circle (0.85pt) node[above right=-1pt] {$C$};
\end{tikzpicture}
\end{center}
\nopagebreak[4]
\Needspace{4\baselineskip}
\textbf{解析}

设侧棱长为 $l$，侧面等腰三角形的底角为 $\theta$（侧面的底边就是底面边长 $2$）。

\Needspace{5\baselineskip}
展开图中 $B$ 落在 $P_1P_2$ 上，$B$ 一侧的三个角拼成平角：
\nopagebreak[4]
\[\angle P_1BA+\angle ABC+\angle CBP_2=180^\circ .\]
\Needspace{5\baselineskip}
正三棱锥的三个侧面全等，故 $\angle P_1BA=\angle CBP_2=\theta$；底面是正三角形，$\angle ABC=60^\circ$。于是
\nopagebreak[4]
\[2\theta+60^\circ=180^\circ,\qquad\theta=60^\circ .\]
\Needspace{5\baselineskip}
每个侧面都是边长为 $2$ 的正三角形，$l=2$，从而展开图的三边为
\nopagebreak[4]
\[P_1P_2=P_2P_3=P_3P_1=2l=4 .\]

\Needspace{5\baselineskip}
此时三棱锥各棱长均为 $2$，是正四面体。底面外接圆半径
\nopagebreak[4]
\[r=\frac{2}{2\sin60^\circ}=\frac{2}{\sqrt3},\]
\Needspace{5\baselineskip}
高
\nopagebreak[4]
\[h=\sqrt{l^2-r^2}=\sqrt{4-\frac43}=\frac{2\sqrt6}{3},\]
\Needspace{5\baselineskip}
故
\nopagebreak[4]
\[V=\frac13\cdot\frac{\sqrt3}{4}\cdot2^2\cdot\frac{2\sqrt6}{3}=\frac{2\sqrt2}{3} .\]

\gkanswer{答案：$P_1P_2=P_2P_3=P_3P_1=4$，$V=\dfrac{2\sqrt2}3$。展开图为三角形迫使侧面为正三角形，侧棱 $l=2$，各棱长全为 $2$。}

\end{answer}
\begin{answer}
\gksolhead{20}{}

\textbf{题目}：

设常数 $a\geqslant0$，函数 $f(x)=\dfrac{2^x+a}{2^x-a}$。

（1） 若 $a=4$，求函数 $y=f(x)$ 的反函数 $y=f^{-1}(x)$；

（2） 根据 $a$ 的不同取值，讨论函数 $y=f(x)$ 的奇偶性，并说明理由。

\Needspace{4\baselineskip}
\textbf{解析}

\gksubpart{（1）}

\Needspace{5\baselineskip}
$a=4$ 时
\nopagebreak[4]
\[f(x)=\frac{2^x+4}{2^x-4},\]
\Needspace{5\baselineskip}
定义域为 $x\neq2$。设 $y=\dfrac{2^x+4}{2^x-4}$，则
\nopagebreak[4]
\[y(2^x-4)=2^x+4\iff(y-1)2^x=4(y+1)\iff2^x=\frac{4(y+1)}{y-1},\]
\Needspace{5\baselineskip}
故
\nopagebreak[4]
\[x=\log_2\frac{4(y+1)}{y-1},\]
\Needspace{5\baselineskip}
其中需 $y\neq1$ 且 $\dfrac{y+1}{y-1}>0$，即 $y<-1$ 或 $y>1$。所以
\nopagebreak[4]
\[f^{-1}(x)=\log_2\frac{4(x+1)}{x-1},\qquad x\in(-\infty,-1)\cup(1,+\infty).\]

\gksubpart{（2）}

\Needspace{5\baselineskip}
对定义域内的 $x$，
\nopagebreak[4]
\[f(-x)=\frac{2^{-x}+a}{2^{-x}-a}=\frac{1+a\cdot2^x}{1-a\cdot2^x}=-\frac{a\cdot2^x+1}{a\cdot2^x-1}.\]

\begin{itemize}
\item 当 $a=0$ 时，$f(x)=\dfrac{2^x}{2^x}=1$，定义域为 $\mathbf R$，$f(-x)=1=f(x)$，是偶函数；
\item 当 $a=1$ 时，定义域为 $\{x\mid x\neq0\}$，关于原点对称，且
\end{itemize}
\nopagebreak[4]
\[f(-x)=-\frac{2^x+1}{2^x-1}=-f(x),\]
是奇函数；
\begin{itemize}
\item 当 $a>0$ 且 $a\neq1$ 时，定义域为 $\{x\mid x\neq\log_2a\}$，因 $\log_2a\neq0$，定义域不关于原点对称，故既不是奇函数也不是偶函数。
\end{itemize}

\gkanswer{答案（1）：$f^{-1}(x)=\log_2\dfrac{4(x+1)}{x-1}$，定义域 $(-\infty,-1)\cup(1,+\infty)$。由 $2^x=\dfrac{4(y+1)}{y-1}$ 解出。}

\gkanswer{答案（2）：$a=0$ 时为偶函数；$a=1$ 时为奇函数；$a>0$ 且 $a\neq1$ 时既非奇函数也非偶函数。$a>0$ 且 $a\neq1$ 时定义域 $\{x\mid x\neq\log_2a\}$ 不关于原点对称。}

\end{answer}
\begin{answer}
\gksolhead{21}{}

\textbf{题目}：

如图，某公司要在 $A$、$B$ 两地连线上的定点 $C$ 处建造广告牌 $CD$，其中 $D$ 为顶端，$AC$ 长 $35$ 米，$CB$ 长 $80$ 米。设点 $A$、$B$ 在同一水平面上，从 $A$ 和 $B$ 看 $D$ 的仰角分别为 $\alpha$ 和 $\beta$。
\nopagebreak[4]
\begin{center}% ============================================================
% f2014_shanghai_arts_q21_billboard.tex
% 2014 上海卷（文科）第 21 题：广告牌 CD 的仰角示意图
% 来源 PDF：普通高考/2014/2014上海文.pdf
% 原页：第 2 页（A3），印面 918，第 21 题题干下方（左栏）
% A、C、B 在同一水平线上，AC=35 米、CB=80 米，D 为广告牌顶端；
% 图中 CD 为铅垂方向（第 (1) 问的情形），α、β 分别为从 A、B
% 看 D 的仰角。示意图按 AC:CB≈35:80 的比例绘制，高度取示意值。
% ============================================================
\begin{tikzpicture}[>=Stealth, line width=0.65pt, font=\small, scale=0.95,
  line cap=round, line join=round]
  \coordinate (A) at (0,0);
  \coordinate (C) at (3.50,0);
  \coordinate (B) at (11.50,0);
  \coordinate (D) at (3.50,3.20);

  % 水平线、广告牌与两条视线
  \draw (A) -- (B);
  \draw (D) -- (C);
  \draw (A) -- (D);
  \draw (B) -- (D);

  % 仰角 α、β 的角标
  \draw (1.05,0) arc (0:42.44:1.05);
  \draw ([shift=(158.13:1.10)]B) arc (158.13:180:1.10);

  % 顶点标注
  \fill (A) circle (0.85pt) node[below left=-1pt] {$A$};
  \fill (C) circle (0.85pt) node[below=-1pt] {$C$};
  \fill (B) circle (0.85pt) node[below right=-1pt] {$B$};
  \fill (D) circle (0.85pt) node[above=-1pt] {$D$};
  \node at (1.60,0.52) {$\alpha$};
  \node at (10.45,0.60) {$\beta$};
\end{tikzpicture}
\end{center}
\nopagebreak[4]
（1） 设计中 $CD$ 是铅垂方向。若要求 $\alpha\geqslant2\beta$，问 $CD$ 的长至多为多少（结果精确到 $0.01$ 米）？

（2） 施工完成后，$CD$ 与铅垂方向有偏差。现在实测得 $\alpha=38.12^\circ$，$\beta=18.45^\circ$，求 $CD$ 的长（结果精确到 $0.01$ 米）。

\Needspace{4\baselineskip}
\textbf{解析}

\gksubpart{（1）}

\Needspace{5\baselineskip}
设 $CD=h$（米）。由 $CD$ 铅垂、$A$，$C$，$B$ 在同一水平线上，得
\nopagebreak[4]
\[\tan\alpha=\frac h{35},\qquad\tan\beta=\frac h{80}.\]
\Needspace{5\baselineskip}
若 $h\geqslant80$，则 $\tan\beta\geqslant1$，即 $2\beta\geqslant\dfrac\pi2>\alpha$，不满足 $\alpha\geqslant2\beta$，故只需讨论 $h<80$，此时 $2\beta<\dfrac\pi2$。由 $\alpha\geqslant2\beta$ 得 $\tan\alpha\geqslant\tan2\beta$，而
\nopagebreak[4]
\[\tan2\beta=\frac{2\tan\beta}{1-\tan^2\beta}=\frac{2\cdot\frac h{80}}{1-\left(\frac h{80}\right)^2}=\frac{160h}{6400-h^2},\]
\Needspace{5\baselineskip}
故
\nopagebreak[4]
\[\frac h{35}\geqslant\frac{160h}{6400-h^2}.\]
\Needspace{5\baselineskip}
因 $h>0$，约去 $h$ 得
\nopagebreak[4]
\[\frac1{35}\geqslant\frac{160}{6400-h^2}\iff6400-h^2\geqslant5600\iff h^2\leqslant800,\]
故 $h\leqslant20\sqrt2\approx28.28$。所以 $CD$ 的长至多为 $20\sqrt2$ 米，约 $28.28$ 米。

\gksubpart{（2）}

\Needspace{5\baselineskip}
设 $D$ 在地面上的投影为 $D'$，$AD'=x$ 米，$DD'=h$ 米，则
\nopagebreak[4]
\[\tan\alpha=\frac hx,\qquad\tan\beta=\frac h{115-x},\qquad x\in(0,115).\]
\Needspace{5\baselineskip}
由 $h=x\tan\alpha=(115-x)\tan\beta$ 得
\nopagebreak[4]
\[x=\frac{115\tan\beta}{\tan\alpha+\tan\beta}.\]
\Needspace{5\baselineskip}
取 $\tan38.12^\circ\approx0.7847$，$\tan18.45^\circ\approx0.3336$，则
\nopagebreak[4]
\[x\approx\frac{115\times0.3336}{0.7847+0.3336}\approx34.31,\qquad h=x\tan\alpha\approx34.31\times0.7847\approx26.92 .\]
\Needspace{5\baselineskip}
于是
\nopagebreak[4]
\[CD=\sqrt{(x-35)^2+h^2}\approx\sqrt{(-0.69)^2+26.92^2}\approx26.93 .\]
故施工完成后 $CD$ 的长约为 $26.93$ 米。

\gkanswer{答案（1）：至多 $20\sqrt2$ 米，约 $28.28$ 米。由 $\tan\alpha\geqslant\tan2\beta$ 得 $6400-h^2\geqslant5600$。}

\gkanswer{答案（2）：约 $26.93$ 米。由 $x=\dfrac{115\tan\beta}{\tan\alpha+\tan\beta}$、$h=x\tan\alpha$ 得水平偏差约 $0.69$ 米、高约 $26.92$ 米，$CD=\sqrt{0.69^2+26.92^2}$。}

\end{answer}
\begin{answer}
\gksolhead{22}{}

\textbf{题目}：

在平面直角坐标系 $xOy$ 中，对于直线 $l:ax+by+c=0$ 和点 $P_1(x_1,y_1)$，$P_2(x_2,y_2)$，记 $\eta=(ax_1+by_1+c)(ax_2+by_2+c)$。若 $\eta<0$，则称点 $P_1$，$P_2$ 被直线 $l$ 分隔。若曲线 $C$ 与直线 $l$ 没有公共点，且曲线 $C$ 上存在点 $P_1$，$P_2$ 被直线 $l$ 分隔，则称直线 $l$ 为曲线 $C$ 的一条分隔线。

（1） 求证：点 $A(1,2)$，$B(-1,0)$ 被直线 $x+y-1=0$ 分隔；

（2） 若直线 $y=kx$ 是曲线 $x^2-4y^2=1$ 的分隔线，求实数 $k$ 的取值范围；

（3） 动点 $M$ 到点 $Q(0,2)$ 的距离与到 $y$ 轴的距离之积为 $1$，设点 $M$ 的轨迹为 $E$。求 $E$ 的方程，并证明 $y$ 轴为曲线 $E$ 的分隔线。

\Needspace{4\baselineskip}
\textbf{解析}

\gksubpart{（1）}

\Needspace{5\baselineskip}
对直线 $x+y-1=0$ 有 $a=b=1$，$c=-1$，于是
\nopagebreak[4]
\[\eta=(1+2-1)\times(-1+0-1)=2\times(-2)=-4<0,\]
故点 $A(1,2)$，$B(-1,0)$ 被直线 $x+y-1=0$ 分隔。

\gksubpart{（2）}

\Needspace{5\baselineskip}
把 $y=kx$ 代入 $x^2-4y^2=1$ 得
\nopagebreak[4]
\[(1-4k^2)x^2=1 .\]
当 $1-4k^2>0$ 即 $|k|<\dfrac12$ 时，$x=\pm\dfrac1{\sqrt{1-4k^2}}$，直线与曲线有公共点；当 $|k|\geqslant\dfrac12$ 时方程无解，直线与曲线没有公共点。

\Needspace{5\baselineskip}
再验证 $|k|\geqslant\dfrac12$ 时曲线上的点确能被分隔。取 $t>1$，令 $u=\dfrac{\sqrt{t^2-1}}2$，则
\nopagebreak[4]
\[P_1(t,u),\qquad P_2(-t,u)\]
\Needspace{5\baselineskip}
都在 $x^2-4y^2=1$ 上，且
\nopagebreak[4]
\[\eta=(kt-u)(-kt-u)=u^2-k^2t^2 .\]
\Needspace{5\baselineskip}
由 $|k|\geqslant\dfrac12$ 得
\nopagebreak[4]
\[u^2-k^2t^2=\frac{t^2-1}4-k^2t^2\leqslant\frac{t^2-1}4-\frac{t^2}4=-\frac14<0,\]
即 $P_1$，$P_2$ 被直线 $y=kx$ 分隔。

所以 $k$ 的取值范围为 $\left(-\infty,-\dfrac12\right]\cup\left[\dfrac12,+\infty\right)$。

\gksubpart{（3）}

\Needspace{5\baselineskip}
点 $M(x,y)$ 到 $Q(0,2)$ 的距离为 $\sqrt{x^2+(y-2)^2}$，到 $y$ 轴的距离为 $|x|$，故
\nopagebreak[4]
\[\sqrt{x^2+(y-2)^2}\cdot|x|=1 .\]
\Needspace{5\baselineskip}
两边平方得 $E$ 的方程
\nopagebreak[4]
\[x^2\left[x^2+(y-2)^2\right]=1 .\]
该方程左端非负、右端为 $1$，故 $x\neq0$，$E$ 与 $y$ 轴没有公共点。

\Needspace{5\baselineskip}
又 $x=\pm1$，$y=2$ 时左端为 $1$，故 $(1,2)$，$(-1,2)$ 都在 $E$ 上。对直线 $x=0$（即 $a=1$，$b=0$，$c=0$），
\nopagebreak[4]
\[\eta=1\times(-1)=-1<0,\]
即这两点被 $y$ 轴分隔。所以 $y$ 轴是曲线 $E$ 的一条分隔线。

\gkanswer{答案（1）：证明。$\eta=(1+2-1)(-1+0-1)=2\times(-2)=-4<0$，满足分隔定义。}

\gkanswer{答案（2）：$k\in\left(-\infty,-\dfrac12\right]\cup\left[\dfrac12,+\infty\right)$。$|k|<\dfrac12$ 时直线与双曲线相交；$|k|\geqslant\dfrac12$ 时无公共点，且取关于 $y$ 轴对称的两点可使 $\eta<0$。}

\gkanswer{答案（3）：$E:x^2\left[x^2+(y-2)^2\right]=1$；证明见解析。$y$ 轴与 $E$ 无公共点，且 $E$ 上两点 $(1,2)$、$(-1,2)$ 被 $y$ 轴分隔。}

\end{answer}
\begin{answer}
\gksolhead{23}{}

\textbf{题目}：

已知数列 $\{a_n\}$ 满足 $\dfrac13a_n\leqslant a_{n+1}\leqslant3a_n$，$n\in\mathbf N^*$，$a_1=1$。

（1） 若 $a_2=2$，$a_3=x$，$a_4=9$，求 $x$ 的取值范围；

（2） 若 $\{a_n\}$ 是等比数列，且 $a_m=\dfrac1{1000}$，求正整数 $m$ 的最小值，以及 $m$ 取最小值时相应 $\{a_n\}$ 的公比；

（3） 若 $a_1,a_2,\cdots,a_{100}$ 成等差数列，求 $a_1,a_2,\cdots,a_{100}$ 的公差的取值范围。

\Needspace{4\baselineskip}
\textbf{解析}

\gksubpart{（1）}

\Needspace{5\baselineskip}
由 $\dfrac13a_2\leqslant a_3\leqslant3a_2$ 得
\nopagebreak[4]
\[\frac23\leqslant x\leqslant6 .\]
\Needspace{5\baselineskip}
由 $\dfrac13a_3\leqslant a_4\leqslant3a_3$ 得
\nopagebreak[4]
\[\frac x3\leqslant9\leqslant3x\iff3\leqslant x\leqslant27 .\]
取交集得 $3\leqslant x\leqslant6$，即 $x$ 的取值范围为 $[3,6]$。

\gksubpart{（2）}

\Needspace{5\baselineskip}
设 $a_n=q^{n-1}$。若 $q<0$，则 $a_1=1>0>a_2=q$，而条件要求 $a_2\geqslant\dfrac13a_1>0$，矛盾；故 $q>0$。约去 $a_n>0$ 得
\nopagebreak[4]
\[\frac13\leqslant q\leqslant3 .\]
\Needspace{5\baselineskip}
又 $a_m=q^{m-1}=\dfrac1{1000}$，故
\nopagebreak[4]
\[q=1000^{-\frac1{m-1}}=10^{-\frac3{m-1}} .\]

\Needspace{5\baselineskip}
当 $2\leqslant m\leqslant7$ 时，$m-1\leqslant6$，于是
\nopagebreak[4]
\[q=10^{-\frac3{m-1}}\leqslant10^{-\frac36}=10^{-\frac12}=\frac1{\sqrt{10}}<\frac13,\]
不满足 $q\geqslant\dfrac13$。

\Needspace{5\baselineskip}
当 $m=8$ 时，
\nopagebreak[4]
\[q=10^{-\frac37},\]
\Needspace{5\baselineskip}
且
\nopagebreak[4]
\[q\geqslant\frac13\iff10^{-\frac37}\geqslant\frac13\iff3^7\geqslant10^3,\]
而 $3^7=2187>1000=10^3$，故 $q\in\left[\dfrac13,3\right]$。

\Needspace{5\baselineskip}
所以 $m$ 的最小值为 $8$，此时公比
\nopagebreak[4]
\[q=10^{-\frac37}=\sqrt[7]{\frac1{1000}}\approx0.373 .\]

\gksubpart{（3）}

\Needspace{5\baselineskip}
设公差为 $d$，则 $a_n=1+(n-1)d$。条件对 $n=1,2,\cdots,99$ 给出
\nopagebreak[4]
\[\frac13a_n\leqslant a_{n+1}\leqslant3a_n .\]

\Needspace{5\baselineskip}
左半式：
\nopagebreak[4]
\[\begin{aligned}
1+nd&\geqslant\frac{1+(n-1)d}3\\
\iff 3+3nd&\geqslant1+(n-1)d\\
\iff 2+(2n+1)d&\geqslant0,
\end{aligned}\]
\Needspace{5\baselineskip}
即
\nopagebreak[4]
\[d\geqslant-\frac2{2n+1}\qquad(n=1,\cdots,99).\]
\Needspace{5\baselineskip}
右端随 $n$ 增大而增大，故只需 $n=99$：
\nopagebreak[4]
\[d\geqslant-\frac2{199} .\]

\Needspace{5\baselineskip}
右半式：
\nopagebreak[4]
\[\begin{aligned}
1+nd&\leqslant3\left[1+(n-1)d\right]\\
\iff 2+(2n-3)d&\geqslant0 .
\end{aligned}\]
\Needspace{5\baselineskip}
当 $n=1$ 时得 $d\leqslant2$；当 $n\geqslant2$ 时得
\nopagebreak[4]
\[d\geqslant-\frac2{2n-3}\qquad(n=2,\cdots,99),\]
\Needspace{5\baselineskip}
右端随 $n$ 增大而增大，故只需 $n=99$：
\nopagebreak[4]
\[d\geqslant-\frac2{195}.\]
而 $-\dfrac2{195}<-\dfrac2{199}$，故这一限制弱于左半式给出的 $d\geqslant-\dfrac2{199}$。

\Needspace{5\baselineskip}
综上，
\nopagebreak[4]
\[-\frac2{199}\leqslant d\leqslant2 .\]

\Needspace{5\baselineskip}
当 $d=-\dfrac2{199}$ 时，$a_{100}=1+99d=\dfrac1{199}>0$，且
\nopagebreak[4]
\[a_{n+1}\geqslant\frac13a_n\iff199\geqslant2n+1\iff n\leqslant99,\]
\nopagebreak[4]
\[a_{n+1}\leqslant3a_n\iff199\geqslant2n-3\iff n\leqslant101,\]
对 $n=1,\cdots,99$ 均成立；$d=2$ 时 $a_n=2n-1$，$a_{n+1}=a_n+2\leqslant3a_n$ 且 $a_{n+1}>a_n>\dfrac13a_n$。故端点可取到。

\Needspace{5\baselineskip}
还需说明 $n=100$ 不妨碍取值：此时只需 $a_{100}>0$，便可取 $a_{101}=a_{100}$ 并使之后各项不变，而由 $d\geqslant-\dfrac2{199}$ 得
\nopagebreak[4]
\[a_{100}=1+99d\geqslant1-\frac{198}{199}=\frac1{199}>0 .\]
\Needspace{5\baselineskip}
故公差的取值范围为
\nopagebreak[4]
\[\left[-\frac2{199},2\right].\]

\gkanswer{答案（1）：$[3,6]$。由 $\dfrac13x\leqslant9\leqslant3x$ 得 $3\leqslant x\leqslant27$，与 $\dfrac23\leqslant x\leqslant6$ 取交。}

\gkanswer{答案（2）：$m$ 的最小值为 $8$，此时公比 $q=10^{-\frac37}$。由 $q^{m-1}=10^{-3}$ 及 $\dfrac13\leqslant q\leqslant3$ 得 $m-1\geqslant\log_31000$。}

\gkanswer{答案（3）：$\left[-\dfrac2{199},2\right]$。由 $\dfrac13a_n\leqslant a_{n+1}\leqslant3a_n$ 得 $d\geqslant-\dfrac2{2n+1}$ 与 $d\leqslant2$（$n=1,\cdots,99$）。}

\end{answer}